\scanpage{48}
\begin{center}2 -- 12 --\end{center}
\begin{flushright}n\textsuperscript{o} 392\end{flushright}

c) Le corps résiduel de $v$ s'identifie à $A/p$, et celui de $v_i$ à $B/p_i$, (\S\ 1, n\textsuperscript{o} 4, cor. 1 de la prop. 6).

d) Soit $a$, (resp. $a_i$), une uniformisante pour $v$, (resp. $v_i$). On a
\[
 aB_{p_i}=aA_pB_{p_i}=pA_pB_{p_i}=pB.B_{p_i}
 =(\prod_j p_j^{e(j)})B_{p_i}
 =\prod_j(p_jB_{p_i})^{e(j)}=(p_iB_{p_i})^{e(i)}=
\]
\[
 =a_i^{e(i)}B_{p_i}\qquad\text{puisque }p_jB_{p_i}=B_{p_i}\quad\text{pour }j\ne i\ ;\ \text{d'où d), puisque}
\]
\[
 e(v_i/v)=v_i(a).
\]
\begin{center}-:-:-:-\end{center}
\mspage{47}
